\chapter{HASIL DAN PEMBAHASAN}
\label{chap:bab4}

% Catatan:
% Bab 4 berfokus pada HASIL implementasi dan PEMBAHASAN dari metode pada Bab 3.
% Hindari mengulang teori/definisi; cukup tunjukkan keluaran (tabel/gambar) dan analisisnya.

\section{Hasil Implementasi}
\label{sec:hasil-implementasi}
% Isi ringkas:
% - Ringkasan alur implementasi end-to-end (data -> preprocessing -> model -> evaluasi)
% - Lingkungan eksperimen: spesifikasi perangkat (CPU/GPU), OS, versi library utama
% - Reproducibility: seed, pembagian data, dan skenario pengujian

\subsection{Analisis Dataset}
\label{subsec:analisis-dataset}
% Letakkan "cerita data" dan bukti kuantitatif sebelum/sesudah trimming di sini.

\subsubsection{Deskripsi dan Distribusi Kelas}
\label{subsubsec:deskripsi-distribusi}
% Isi:
% - Jumlah kelas gesture
% - Definisi singkat tiap kelas (opsional 1 kalimat/kode label)
% - Distribusi jumlah sampel per kelas (setelah trimming)
% - Tampilkan tabel distribusi dan/atau bar chart

\subsubsection{\textit{Trimming} Data}
\label{subsubsec:sebelum-sesudah-trimming}
% Isi:
% - Jelaskan dataset mentah (jumlah video awal, durasi rata-rata, kondisi rekaman)
% - Jelaskan hasil setelah trimming/segmentasi (jumlah segmen/sampel menjadi N)
% - Tabel WAJIB: sebelum vs sesudah trimming (jumlah video/segmen, durasi, total frame)
% - Tekankan alasan trimming dilakukan (mengurangi bagian non-gesture, meningkatkan fokus sinyal)

\subsubsection{Contoh Visual Sampel Dataset}
\label{subsubsec:contoh-visual-dataset}
% Isi:
% - Gambar/screenshot representatif:
%   (1) frame sebelum trimming (scene penuh)
%   (2) frame sesudah trimming (gesture lebih fokus)
% - Jika memungkinkan tambahkan 1 contoh per beberapa kelas (secukupnya, jangan terlalu banyak)

\subsection{Hasil Prapemrosesan}
\label{subsec:hasil-prapemrosesan}
% Prinsip: tampilkan OUTPUT preprocessing (bukan langkah teori).
% Sertakan contoh sebelum/sesudah, dan ringkas parameter penting.

\subsubsection{Hasil Segmentasi dan Trimming}
\label{subsubsec:hasil-segmentasi-trimming}
% Isi:
% - Visual: sebelum vs sesudah trimming (minimal 1-2 contoh)
% - Tabel opsional: statistik durasi/frame sebelum vs sesudah
% - Jelaskan kriteria potong (misal start/end gesture) dan dampak ke jumlah sampel

\subsubsection{Hasil Ekstraksi Landmark}
\label{subsubsec:hasil-ekstraksi-landmark}
% Isi:
% - Gambar WAJIB: frame dengan overlay landmark (misal MediaPipe)
% - Jelaskan landmark apa yang dipakai (hands/pose/face) secara ringkas
% - Jelaskan kualitas deteksi (misal kasus gagal deteksi, noise) + cara penanganan (jika ada)

\subsubsection{Seleksi Titik Landmark dan Normalisasi}
\label{subsubsec:seleksi-normalisasi}
% Isi:
% - Titik mana yang dipakai (misal hanya tangan, atau subset)
% - Teknik normalisasi (misal relative coordinate, scaling)
% - Alasan: mengurangi variansi posisi kamera/jarak subjek

\subsubsection{Fitur Temporal dan Representasi Urutan}
\label{subsubsec:fitur-temporal}
% Isi:
% - Jelaskan bentuk fitur final yang dipakai untuk eksperimen:
%   (a) koordinat (x,y)
%   (b) koordinat + delta (x,y,Δx,Δy)
% - Visual opsional: plot perubahan koordinat terhadap waktu untuk 1 gesture
% - Kaitkan secara singkat ke tujuan studi ablasi (pengaruh fitur temporal)

\subsubsection{Bentuk Masukan Model (Tensor Input)}
\label{subsubsec:tensor-input}
% Isi:
% - Tabel WAJIB: dimensi input (C x T x L) atau format lain yang digunakan
% - Jelaskan arti C (channel), T (time/frame), L (jumlah landmark/titik)
% - Jelaskan perbedaan skenario input 2-channel vs 4-channel (untuk xy vs xy+dxdy)

\subsection{Implementasi Model Klasifikasi}
\label{subsec:implementasi-model}
% Isi:
% - Arsitektur yang digunakan (ResNet18/34/50)
% - Adaptasi layer input (conv1) untuk jumlah channel yang berbeda
% - Skema training: loss, optimizer, learning rate, epoch, batch size
% - Regularisasi/augmentasi (jika ada)
% - Tabel WAJIB: hyperparameter dan konfigurasi training
% - (Opsional) Singgung waktu training dan penggunaan GPU

\subsubsection{Konfigurasi dan Modifikasi Arsitektur}
\label{subsubsec:konfigurasi-arsitektur}
% Isi:
% - Bagian mana yang dimodifikasi dari ResNet standar
% - Alasan pemilihan ResNet (stabil, umum, mudah dibandingkan)
% - Skenario: input 2 channel vs 4 channel

\subsubsection{Pengaturan Pelatihan dan Pembagian Data}
\label{subsubsec:pengaturan-pelatihan}
% Isi:
% - Train/val/test split atau cross-validation (jelaskan)
% - Kriteria early stopping / checkpoint (jika ada)
% - Seed dan cara memastikan eksperimen fair antar model

\section{Evaluasi Hasil}
\label{sec:evaluasi-hasil}
% Isi:
% - Definisi metrik yang dipakai (accuracy, precision, recall, F1) secara singkat (1-2 kalimat)
% - Jelaskan protokol evaluasi (test set tetap, rata-rata beberapa run, dsb.)

\subsection{Hasil Baseline Model}
\label{subsec:hasil-baseline}
% Isi:
% - Tabel WAJIB: performa baseline (misal ResNet18 dengan fitur dasar tertentu)
% - Jika ada beberapa baseline, urutkan dari paling sederhana
% - Jelaskan hasil utama secara singkat: mana yang paling baik pada baseline dan indikasi error

\subsubsection{Akurasi dan F1-Score}
\label{subsubsec:akurasi-f1}
% Isi:
% - Tabel metrik utama
% - (Opsional) sertakan std dev jika multiple run

\subsubsection{Confusion Matrix}
\label{subsubsec:confusion-matrix}
% Isi:
% - GAMBAR WAJIB: confusion matrix untuk model terbaik (atau baseline vs terbaik)
% - Bahas kelas mana yang sering tertukar dan kemungkinan penyebab (kemiripan gesture, noise landmark)

\subsection{Studi Ablasi}
\label{subsec:studi-ablasi}
% Ini inti kontribusi: menjawab pertanyaan riset melalui perbandingan terkontrol.

\subsubsection{Pengaruh Penambahan Fitur Temporal}
\label{subsubsec:ablasi-fitur-temporal}
% Isi:
% - Tabel WAJIB: bandingkan xy vs xy+dxdy pada arsitektur yang sama
% - Analisis: apakah naik/turun, pada kelas mana paling membantu (rujuk confusion matrix jika relevan)
% - Kesimpulan mini: apakah Δx,Δy efektif

\subsubsection{Pengaruh Kedalaman Arsitektur ResNet}
\label{subsubsec:ablasi-kedalaman}
% Isi:
% - Tabel WAJIB: ResNet18 vs ResNet34 vs ResNet50 pada fitur yang sama
% - Analisis: indikasi overfitting/underfitting (gap train-val, performa test)
% - Kesimpulan mini: kedalaman optimal untuk dataset ini

\subsubsection{Interaksi Fitur dan Arsitektur}
\label{subsubsec:ablasi-interaksi}
% Isi:
% - Analisis gabungan: apakah Δx,Δy lebih berguna pada ResNet tertentu
% - Sajikan tabel ringkas (model x fitur) dan sebutkan kombinasi terbaik
% - Tarik implikasi praktis (kompleksitas vs performa)

\section{Pembahasan}
\label{sec:pembahasan}
% Isi:
% - Interpretasi hasil utama: mengapa kombinasi terbaik bekerja
% - Kaitan ke teori di Bab 2 (misal pentingnya dinamika gerak untuk gesture)
% - Keterbatasan:
%   * kualitas landmark (occlusion, blur)
%   * jumlah data per kelas (imbalanced)
%   * generalisasi ke subjek baru/latar berbeda
% - Implikasi: apa yang bisa ditingkatkan ke depan (augmentasi, model temporal, dsb.)

\section{Rangkuman Hasil}
\label{sec:rangkuman-hasil}
% Isi:
% - Poin-poin ringkas (bullet) dari:
%   * hasil dataset (sebelum/sesudah trimming)
%   * hasil preprocessing (output final tensor)
%   * model terbaik + angka metrik utama
%   * temuan ablation (Δx,Δy dan kedalaman)
% - Menjadi jembatan menuju Bab 5 (kesimpulan dan saran)
% Isi:
% - Jawab RM/RQ secara naratif berbasis angka:
%   RQ1: pengaruh Δx,Δy -> rujuk subsubsec:ablasi-fitur-temporal
%   RQ2: pengaruh kedalaman ResNet -> rujuk subsubsec:ablasi-kedalaman
% - Jangan buat "Evaluasi RM" terpisah; cukup simpulkan dari hasil ablasi
